\section{Introduction}
\label{sec:intro}

Modern network intrusion detection systems (NIDS) increasingly rely on machine learning over flow-based telemetry. At multi-gigabit rates, inspecting every packet payload is prohibitive, so sensors such as CICFlowMeter \cite{sharafaldin2018toward} aggregate bidirectional packet exchanges into tabular flow records of timing, length, header and flag statistics. Gradient-boosted decision trees (GBDTs), notably XGBoost \cite{chen2016xgboost}, LightGBM \cite{ke2017lightgbm} and CatBoost \cite{prokhorenkova2018catboost}, are the dominant classifiers for these records \cite{dang2022using, dang2022studying}.

Near-perfect benchmark accuracy does not guarantee robustness: attackers control the malleable parts of their traffic, such as timing, padding, segmentation and window sizes, so a detector that relies on these features can be evaded without losing the exploit \cite{apruzzese2022modeling, sommer2010outside, geirhos2020shortcut}.

Evading a tree ensemble is hard for two reasons. First, the ensemble is piecewise constant, so $\nabla_x F(x)=0$ for almost all $x$ and gradient-based attacks receive no signal. Second, flow features obey conservation laws. Flow Bytes/s satisfies $x_8 = 10^6(x_4+x_5)/x_6$, where $x_4$ and $x_5$ are the forward and backward byte totals and $x_6$ is the duration in microseconds. A perturbation that changes the byte totals without updating the rate produces a phantom flow that protocol sanity checks reject \cite{pierazzi2020intrusion, apruzzese2022modeling}. Existing attacks either ignore this coupling, as ZOO \cite{chen2017zoo} and HopSkipJump \cite{chen2020hopskipjump} do (HopSkipJump yields no feasible evasive flow on the corpora of this study; ZOO keeps at most 25.2\% Real ASR, the share of attacks that evade and satisfy the schema invariants), or need hundreds of black-box queries per flow \cite{simonetto2024constrained, ghamizi2020moeva, han2021traffic}, or depend on a surrogate \cite{simonetto2024constrained, zolbayar2022nidsgan}.

Our earlier attack, Flow-Field Evasion (FFE-v1), derived four packet-level actions as closed-form generators of a Lie cone on the flow manifold, so that every step preserves the invariants, and steered with score probes. It evaded XGBoost in 99.5\%, 99.2\% and 57.6\% of the detected attacks of three full corpora, but a rigorous control exposed its limits: a greedy white-box attack in the style of leaf-tuple search \cite{zhang2020efficient}, which reads the trees but projects every candidate onto the manifold, reached 90.2\% on UNSW-NB15, and FFE-v1 evaded LightGBM on CIC-IDS2017 in only 35.1\% and CatBoost on CSE-CIC-IDS2018 in only 59.3\% of the flows. We diagnosed these gaps on a validation split, not on the test cohorts: the attack lacked a generator that raises the packet count without adding payload, used one fixed magnitude per generator and one hard-coded window target (64,240~B) that lies outside the support of the CIC-IDS2017 data, and measured progress on a saturating probability.

This article reports the repaired attack, FFE-v2, and a defense built on what the attack revealed. The main contributions are as follows.
\begin{itemize}
    \item \textbf{Generators.} We extend the Lie-cone model by a re-segmentation generator (the same payload carried by more, smaller packets). The five fields still commute and preserve the rate, mean-length, header and MTU invariants of the schema, although they are not linearly independent (re-segmentation lies in the span of chaff and padding, outside their nonnegative cone; Section~\ref{sec:framework}, Online Resource~1, Section~S1).
    \item \textbf{Attack.} FFE-v2 probes every generator on a geometric ladder of magnitudes, draws window targets from the attacker's benign sample, measures progress in log-odds and refines the cheapest crossing; it needs only scores (Section~\ref{sec:algorithm}). In a single evaluation on frozen settings it matches or exceeds our greedy threshold-crossing white-box control in all nine corpus-library cells (significantly above it in four) and is significantly above FFE-v1 in seven (Real ASR 99.5\%, 100.0\% and 98.3\% on XGBoost against 99.5\%, 99.2\% and 57.6\% for FFE-v1 and 88.5\%, 89.5\% and 21.1\% for CAPGD); five-seed means agree (Section~\ref{sec:results}). The control is our own implementation (erratic across seeds on CIC-IDS2017 LightGBM); LT-Attack and Cascioli et al.\ were not run.
    \item \textbf{Honest accounting.} FFE-v2 uses 28--38 score queries (109 on CIC-IDS2017 CatBoost), against 2 for the control and about 3 for CAPGD with an uncounted surrogate, and more than FFE-v1 on the easy corpora; its operational cost on UNSW-NB15 exceeds that of FFE-v1; one cell (CatBoost on CIC-IDS2017, 36.4\%) is a ceiling set by attacker-immutable features rather than an attack success. ``Feasible'' is relative to a schema verifier that does not check the header-plus-payload bound, and chaff probes are booked as $L_0$ bytes of the length column; a synthetic replay shows that UNSW-NB15 evasions depend on chaff accounting (Section~\ref{sec:ablation}).
    \item \textbf{Defense.} We OR-fuse a monotone gate, a tree ensemble that is non-decreasing in forward packets, forward bytes, duration and header bytes and does not read the window, rates or mean length, with a free classifier (augmentation-hardened on five of nine cells; the undefended model on the others). If the attacker leaves the backward features and the port and protocol unchanged (an assumption on which the guarantee rests), a flow the gate flags stays flagged (Section~\ref{sec:defense}, a certificate in the spirit of monotone classification \cite{incer2018adversarially}). Against FFE-v1 ($M=5,20$), CAPGD and a greedy white-box threshold-crossing attacker, Real ASR on the flows still detected falls to at most 8.8\% on CSE-CIC-IDS2018 and CIC-IDS2017 and to 1.0--32\% on UNSW-NB15, with validation accuracy, F1 and recall not falling by more than 0.5 point. We state where this fails: about 8\% of CSE-CIC-IDS2018 flows are not flagged by the gate and remain evadable, the stand-alone gate costs accuracy, hiding the backward features from the gate to remove the fixed-backward assumption costs much accuracy and weakens V2, and FFE-v2 was run against the hardened models on one cell only.
\end{itemize}
Caveat: FFE-v2 was tuned on validation cohorts and evaluated once on the test cohorts, but the FFE-v1 failures that motivated it were first seen on those test cohorts, re-segmentation was motivated by the control's moves on the same corpora, and the window pool uses the attacker's benign training sample (tuning queries are not counted); the five-seed study (retrained models, fresh cohorts, no tuning) is the cleaner confirmation. The attacks are red teaming that estimates what a schema-constrained adversary with score access can reach; no real network was attacked.
